% from ega4-16-fr.tex
\oldpage[IV-4]{7}
\item[(ii)] Les $Y^{(n)}$ sont des préschémas, canoniquement isomorphes à des
sous-préschémas de $X$.
\item[(iii)] Tout homomorphisme de faisceaux d'anneaux
$\lambda_n:\mathcal{O}_Y\to\mathcal{O}_{Y^{(n)}}$, inverse à droite de
l'homomorphisme d'augmentation $\varphi_{0n}$, fait de
$\mathcal{O}_{Y^{(n)}}$ et des $\mathcal{O}_{Y^{(k)}}$ pour $k\leq n$, des
$\mathcal{O}_Y$-Algèbres quasi-cohérentes; les structures de
$\mathcal{O}_Y$-Module déduites des structures précédentes sur les
$\mathscr{G}r_k(f)$ pour $k\leq n$ coïncident avec celles définies dans
(16.1.2).
\end{enumerate}
\end{proposition}

\noindent(i) La question étant locale sur $X$ et sur $Y$, on peut se borner au cas
où $Y$ est un sous-préschéma fermé de $X$ défini par un Idéal quasi-cohérent
$\mathscr{I}$ de $\mathcal{O}_X$; comme $\mathcal{O}_Y$ est la restriction à
$Y$ de $\mathcal{O}_X/\mathscr{I}$, l'assertion (i) est évidente, et $Y^{(n)}$
est le sous-préschéma fermé de $X$ défini par l'Idéal quasi-cohérent
$\mathscr{I}^{n+1}$ de $\mathcal{O}_X$. Enfin, pour prouver (iii), notons que la
donnée de $\lambda_n$ fait de l'Idéal
$\mathscr{I}/\mathscr{I}^{n+1}$ de l'augmentation $\varphi_{0n}$ et de ses
quotients $\mathscr{I}/\mathscr{I}^{k+1}$ ($1\leq k\leq n$) des
$\mathcal{O}_Y$-Modules, et il suffit de prouver par récurrence sur $k$ que les
$\mathscr{I}/\mathscr{I}^{k+1}$ sont des $\mathcal{O}_Y$-Modules
quasi-cohérents et que la structure de $\mathcal{O}_Y$-Module quotient qu'on en
déduit sur $\mathscr{I}^k/\mathscr{I}^{k+1}$ est la même que celle définie en
(16.1.2). La seconde assertion est immédiate,
$\mathscr{I}^k/\mathscr{I}^{k+1}$ étant annulé par
$\mathscr{I}/\mathscr{I}^{n+1}$; la première résulte, par récurrence sur $k$,
de ce qu'elle est triviale pour $k=1$ et de ce que
$\mathscr{I}/\mathscr{I}^{k+1}$ est une extension de
$\mathscr{I}/\mathscr{I}^k$ par
$\mathscr{I}^k/\mathscr{I}^{k+1}$ (III, 1.4.17).

\begin{corollary}[16.1.6]
\label{IV.16.1.6}
Sous les hypothèses générales de (16.1.5), si l'immersion $f$ est localement de
présentation finie, les $\mathscr{G}r_n(f)$ sont des $\mathcal{O}_Y$-Modules
quasi-cohérents de type fini.
\end{corollary}

En effet, avec les notations de la démonstration de (16.1.5), $\mathscr{I}$ est un
Idéal de type fini de $\mathcal{O}_X$ (1.4.7), donc les
$\mathscr{I}^n/\mathscr{I}^{n+1}$ sont des $\mathcal{O}_Y$-Modules de type fini,
d'où la conclusion.

\begin{corollary}[16.1.7]
\label{IV.16.1.7}
Sous les hypothèses générales de (16.1.5), soit $g:X\to Y$ un morphisme de
préschémas inverse à gauche de $f$. Alors, pour tout $n$, le morphisme composé
$(1,\lambda_n):Y^{(n)}\xrightarrow{h_n}X\xrightarrow{g}Y$ définit un
homomorphisme de faisceaux d'anneaux
$\lambda_n:\mathcal{O}_Y\to\mathcal{O}_{Y^{(n)}}$ inverse à droite de
l'augmentation $\varphi_{0n}$, faisant de $\mathcal{O}_{Y^{(n)}}$ une
$\mathcal{O}_Y$-Algèbre quasi-cohérente; pour ces homomorphismes, les
homomorphismes de transition
$\varphi_{nm}:\mathcal{O}_{Y^{(m)}}\to\mathcal{O}_{Y^{(n)}}$ ($n\leq m$) sont
des homomorphismes de $\mathcal{O}_Y$-Algèbres. En outre, si $g$ est localement
de type fini, les $\mathcal{O}_{Y^{(n)}}$ sont des $\mathcal{O}_Y$-Modules
quasi-cohérents de type fini.
\end{corollary}

La première assertion résulte aussitôt des définitions et de (16.1.5). D'autre
part, si $g$ est localement de type fini, alors $f$ est localement de
présentation finie (1.4.3, (v)); les $\mathscr{G}r_n(f)$ étant alors des
$\mathcal{O}_Y$-Modules quasi-cohérents de type fini par (16.1.6), il en est de
même des $\mathcal{O}_Y$-Modules
$\mathscr{I}/\mathscr{I}^{n+1}$, qui sont des extensions d'un nombre fini de
$\mathscr{G}r_k(f)$ (III, 1.4.17).

\begin{proposition}[16.1.8]
\label{IV.16.1.8}
Soient $X$ un préschéma localement noethérien, $j:Y\to X$ une immersion. Alors
les $Y^{(n)}$ sont des préschémas localement noethériens, les
$\mathscr{G}r_n(j)$ sont des $\mathcal{O}_Y$-Modules cohérents, et
$\mathscr{G}r_\bullet(j)$ est un faisceau d'anneaux cohérent sur l'espace $Y$.
\end{proposition}

Tout étant local sur $X$ et $Y$, on est ramené au cas où $X$ est affine et $j$
une immersion fermée, et alors toutes les assertions sont évidentes sauf la
dernière, qui résulte de ce que si $A$ est un anneau noethérien et
$\mathfrak{I}$ un idéal de $A$, $\operatorname{gr}_{\mathfrak{I}}^\bullet(A)$
est un anneau noethérien, compte tenu de l'exactitude du foncteur $\psi^*$ et de
(0\textsubscript{I}, 5.3.7).

\begin{proposition}[16.1.9]
\label{IV.16.1.9}
